\begin{statement}[C.5.1; Known, with derivation: no smooth CM-torsion parameter]
There is no $t\in\mathbb C\setminus\{0,1\}$ for which $E_t$ has
complex multiplication and $P_t$ is torsion. On the smooth models
at $t=0,1$ after base change of orders $3,4$, respectively, the
limiting points have orders $3,2$. Thus these are the only CM
torsion points on the completed family.
\end{statement}
\begin{proof}
The precise torsion theorem used is \cite{BourdonClarkStankewicz2017},
Theorem 4.15(a), on page 23 of arXiv:1411.2742v3:
for an elliptic curve $E/F$ with CM by any imaginary quadratic order,
\[
 F=\mathbb Q(j(E))\quad\Longrightarrow\quad
 E(F)_{\rm tors}\in
 \{0,\mathbb Z/2,\mathbb Z/3,\mathbb Z/4,
       \mathbb Z/6,(\mathbb Z/2)^2\}
\]
up to group isomorphism. The assertion allows nonmaximal orders.
The exceptional restrictions at discriminants $-3,-4$ belong only
to part (b), over the larger field containing the CM field; they
are not hypotheses of part (a).

For the present twist $F=\mathbb Q(t)=\mathbb Q(j)$ because
$j=1728t$. All Weierstrass coefficients and both coordinates of
$P$ belong to $F$, and $\Delta\ne0$ on the smooth base. Therefore
the theorem applies to the actual twist and the actual section,
not just to an isomorphic complex curve. The point $P$ is never
$O$, and $\psi_2(P),\psi_3(P)$ are nonzero there. The remaining
possible orders are $4$ and $6$, whose parameters were computed
exactly above. For order four, $j=55296/5$ is not an algebraic
integer. For order six the two $j$-values have monic minimal
polynomial
\[
 X^2-\frac{10395648}{11}X+\frac{24461180928}{11};
\]
it is irreducible over $\mathbb Q$, and its coefficients are not
integers. Neither root is an algebraic integer. By the integrality
theorem for singular moduli (for example, Theorem 20.12 and
Corollary 20.13 of \cite{Sutherland2023CM}, proved there for all
imaginary quadratic orders), this excludes complex multiplication
in every remaining case. The required hypothesis is simply that
the endomorphism ring is such an order, which is exactly the CM
assumption.

At zero put $t=s^3$, $x=s^4X$, $y=s^6Y$. The smooth equation is
\[
 Y^2=X^3-\tfrac34s(s^3-1)X+\tfrac14(s^3-1)^2,\qquad
 P_s=(0,(s^3-1)/2).
\]
Its central curve is $Y^2=X^3+1/4$ with endomorphism ring
$\mathbb Z[\rho]$, and the limiting point $(0,-1/2)$ has order
three: its horizontal tangent has triple contact, so $2P=-P$.
At one put $t=1+s^4$, $x=s^2X$, $y=s^3Y$. Then
\[
 Y^2=X^3-\tfrac34(1+s^4)^3X+
       \tfrac14(1+s^4)^4s^2,\qquad
 P_s=(0,(1+s^4)^2s/2).
\]
The central curve is $Y^2=X^3-3X/4$, with endomorphism ring
$\mathbb Z[i]$, and its limiting point $(0,0)$ has exact order two.
The displayed equations have nonzero central discriminants; hence
these are genuine smooth models, and their norm limits are those
already evaluated in C.1.2. This proves Lead 11, including its
twist and boundary qualifications.
\end{proof}
